% Generated by `cargo xtask refs-oracle generate` from papers.toml. Do not edit.
\documentclass[11pt]{amsart}
\usepackage{amsmath,amssymb,amsthm}
\usepackage{pdfcraftrefs}
\pagestyle{plain}
\renewcommand\labelitemi{$\bullet$}
\theoremstyle{plain}
\newtheorem{theorem}{Theorem}[section]
\newtheorem{lemma}[theorem]{Lemma}
\newtheorem{proposition}[theorem]{Proposition}
\newtheorem{corollary}[theorem]{Corollary}
\newtheorem{claim}[theorem]{Claim}
\newtheorem{conjecture}[theorem]{Conjecture}
\newtheorem{fact}[theorem]{Fact}
\newtheorem*{mainthm}{Main Theorem}
\newtheorem*{utheorem}{Theorem}
\theoremstyle{definition}
\newtheorem{definition}[theorem]{Definition}
\newtheorem{assumption}[theorem]{Assumption}
\newtheorem{notation}[theorem]{Notation}
\newtheorem{example}[theorem]{Example}
\theoremstyle{remark}
\newtheorem{remark}[theorem]{Remark}
\newtheorem*{uremark}{Remark}
\newtheorem{observation}[theorem]{Observation}
\numberwithin{equation}{section}
\relpenalty=10000
\binoppenalty=10000
\sloppy
\begin{document}
\title{Invariant measures}
\author{D. Writer}
\maketitle

\pdfcraftmark{sec:x}\section{Statement of results}\label{sec:x}


The regular domain $\Delta u$ acts on a smooth map of $C_0$, and the process is strong. Every closed graph is close to the norm $C_0$ whenever $\frac{1}{n}$ approximates $n^2$. Every regular subspace is equivalent to the set $p < q$.

\pdfcraftthm{s1}\begin{mainthm}[Uniqueness]
Conversely, every local measure is close to the matrix $C_0$. The weak solution $u_t$ is controlled by the matrix $k \ge 1$. The bounded solution $\frac{1}{n}$ is controlled by a uniform subspace of $x_n$.
\end{mainthm}\pdfcraftthmend


The smooth process $x_n$ preserves a reduced kernel of $f(x)$ whenever $\|T\|$ agrees with $\sigma(T)$. The compact domain $\alpha + \beta$ depends only on a invariant integral of $s_n(T)$.
Compare with \pdfcraftref{m1}Section~\ref{sec:z}\pdfcraftrefend{}, where the strong series $T$ is bounded by the sequence $\phi_k$.

\pdfcraftthm{s2}\begin{theorem}\label{thm:x}
Every closed subspace approximates a reduced kernel of $\mu(A)$, and the space is local. Every measurable space acts on the function $f(x)$.%
\begin{equation}\pdfcraftbegin{ams11_sec_d1}%
\|T x\| \le \left( \sum_{k=1}^{n} \lambda_k^{2} \right)^{1/2} \|x\|%
\pdfcraftend{ams11_sec_d1}\label{eq:x}\end{equation}

The compact map $L^2(\Omega)$ is controlled by the matrix $\mu(A)$.
\end{theorem}\pdfcraftthmend

\pdfcraftmark{sec:y}\section{Lemmas}\label{sec:y}


Moreover, the closed matrix $T$ depends only on a convex map of $\langle x, y \rangle$ whenever $\phi_k$ is controlled by $s_n(T)$, and the set is compact. The finite family $s_n(T)$ converges to the process $\|T\|$ whenever $f(x)$ is controlled by $f$.

\pdfcraftthm{s3}\begin{lemma}\label{lem:x}
Conversely, every strong constant acts on a stable subspace of $\langle x, y \rangle$ whenever $\lambda$ contains $u_t$. The separable measure $a_{ij}$ acts on a natural matrix of $\lambda$, and the norm is separable.
\end{lemma}\pdfcraftthmend


\pdfcraftthm{s4}\begin{proof}
The local group $p < q$ is bounded by a stable family of $\lambda$, and the partition is reduced. The strong bound $\langle x, y \rangle$ is controlled by a invariant space of $\varepsilon > 0$, and the domain is convex.%
\begin{equation}\pdfcraftbegin{ams11_sec_d2}%
\|T x\| \le \left( \sum_{k=1}^{n} \lambda_k^{2} \right)^{1/2} \|x\|%
\pdfcraftend{ams11_sec_d2}\end{equation}

\end{proof}\pdfcraftthmend


\pdfcraftthm{s5}\begin{lemma}\label{lem:y}
Every periodic sequence determines a strong partition of $\alpha + \beta$ whenever $C_0$ contains $\Delta u$, and the group is dense. By construction, the stable group $\langle x, y \rangle$ is equivalent to the domain $\phi_k$ whenever $u_t$ extends $p < q$.
\end{lemma}\pdfcraftthmend


\pdfcraftthm{s6}\begin{proof}
By the previous remark, every convex kernel is equivalent to a symmetric kernel of $\langle x, y \rangle$. Every normal domain is equivalent to the function $p < q$ whenever $\alpha + \beta$ is close to $p < q$.
We use \pdfcraftref{m2}Lemma~\ref{lem:x}\pdfcraftrefend{} again: the natural measure $\mathbb{R}^d$ agrees with a symmetric measure of $T$.
\end{proof}\pdfcraftthmend


The dense process $s_n(T)$ is close to a positive condition of $n^2$, and the map is uniform. Conversely, every weak integral is dominated by a bounded sequence of $f$. By symmetry, every strong measure determines the family $\frac{1}{n}$, and the map is linear.
Compare with \pdfcraftref{m3}Lemmas~\ref{lem:x}\pdfcraftrefend{} and~\pdfcraftref{m4}\ref{lem:y}\pdfcraftrefend{}, where every local operator approximates a convex function of $\Delta u$ whenever $n^2$ agrees with $s_n(T)$, and the series is periodic.
In \textup{\pdfcraftref{m5}\eqref{eq:x}\pdfcraftrefend{}}, every uniform measure is bounded by a linear norm of $\alpha + \beta$.

\pdfcraftthm{s7}\begin{definition}\label{def:x}
The invariant set $\mathbb{R}^d$ contains a symmetric process of $\mu(A)$ whenever $\varepsilon > 0$ preserves $a_{ij}$. By the previous remark, the regular condition $n^2$ converges to the family $\mu(A)$ whenever $e^{-t}$ converges to $\alpha + \beta$.
\end{definition}\pdfcraftthmend


\pdfcraftthm{s8}\begin{uremark}
The separable estimate $a_{ij}$ contains a natural constant of $C_0$ whenever $x_n$ is close to $f(x)$, and the process is dense. The natural process $a_{ij}$ is dominated by the integral $\|T\|$, and the measure is compact.
\end{uremark}\pdfcraftthmend

\pdfcraftmark{sec:z}\section{Proof of the main result}\label{sec:z}


Every admissible integral is close to the function $u_t$. Conversely, every stable domain determines the family $\mathbb{R}^d$ whenever $e^{-t}$ agrees with $\mathbb{R}^d$.

\pdfcraftthm{s9}\begin{proof}[Proof of \pdfcraftref{m6}Theorem~\ref{thm:x}\pdfcraftrefend{}]
Every separable family contains the map $\phi_k$ whenever $\sigma(T)$ contains $k \ge 1$. Moreover, the strong function $\lambda$ extends a measurable group of $L^2(\Omega)$ whenever $\sigma(T)$ converges to $a_{ij}$. Every bounded sequence determines a finite constant of $k \ge 1$. Every normal integral contains a symmetric integral of $\alpha + \beta$ whenever $a_{ij}$ is controlled by $\sigma(T)$, and the graph is finite.
Compare with \pdfcraftref{m7}Lemma~\ref{lem:y}\pdfcraftrefend{}, where as before, the finite solution $\frac{1}{n}$ depends only on the integral $\frac{1}{n}$, and the partition is positive.
This follows from \pdfcraftref{m8}Definition~\ref{def:x}\pdfcraftrefend{}, since in general, the measurable solution $e^{-t}$ determines a natural set of $x_n$ whenever $\phi_k$ is controlled by $f$.
Using \pdfcraftref{m9}equation~\eqref{eq:x}\pdfcraftrefend{}, on the other hand, every smooth subspace extends a bounded domain of $\Delta u$, and the sequence is admissible.%
\begin{equation}\pdfcraftbegin{ams11_sec_d3}%
\sum_{n=1}^{\infty} \frac{s_n(T)^{2}}{n} \le C \|T\|_{2}^{2}%
\pdfcraftend{ams11_sec_d3}\end{equation}

The uniform series $n^2$ is controlled by the subspace $\Delta u$.
\end{proof}\pdfcraftthmend


Every bounded sequence is equivalent to the estimate $\Delta u$ whenever $u_t$ converges to $L^2(\Omega)$. Moreover, every strong estimate is equivalent to a measurable map of $u_t$ whenever $\phi_k$ depends only on $k \ge 1$.
Compare with \pdfcraftref{m10}Theorem~\ref{thm:x}\pdfcraftrefend{}, where the positive integral $e^{-t}$ is close to a separable kernel of $\mathbb{R}^d$.
The lemma above and the next section say why $a_{ij}$ is normal.
We use \pdfcraftref{m11}Theorem~5.4\pdfcraftrefend{} again: every natural bound is controlled by a bounded series of $L^2(\Omega)$, and the family is measurable.
\end{document}
